\documentclass[11pt]{article}
\usepackage[T1]{fontenc}
\usepackage{textcomp}
\usepackage[margin=0.75in]{geometry}
\usepackage{amsmath,amssymb,amsthm}
\usepackage{array,booktabs}
\usepackage{hyperref}
\setlength{\parindent}{0pt}
\newtheorem{theorem}{Theorem}
\theoremstyle{definition}
\newtheorem{definition}{Definition}
\title{Flow Proof Artifact: prop-18-square-equality-implies-proportional}
\author{Generated by \texttt{flow doc proof}}
\date{Flow Proof Book}
\begin{document}
\maketitle
If the rectangle contained by the first and third equals the square on the second, the three straight lines are proportional.\\[0.5em]
\textbf{Source.} euclid — Elements, Book VI, Proposition 18\\[0.5em]
\bigskip
\noindent{\large \textbf{Derived fact 1.} \textit{If the rectangle contained by the first and third equals the square on the second, the three straight lines are proportional} \hfill the Euclidean plane \cdot Euclid Book VI \cdot Proposition 18: rectangle on extremes equal to square on mean implies three lines proportional}\par
\smallskip\noindent\textit{Coordinate: the Euclidean plane \cdot Euclid Book VI \cdot Proposition 18: rectangle on extremes equal to square on mean implies three lines proportional}\par
\smallskip\noindent\textit{Source: euclid — Elements, Book VI, Proposition 18}\par
\smallskip\noindent\textit{Built on: proposition 17: three proportional lines give rectangle on extremes equal to square on mean, for Euclid Book VI on the Euclidean plane}\par
\medskip
\noindent\textbf{Goal.} If the rectangle contained by the first and third equals the square on the second, the three straight lines are proportional\par
\noindent$\displaystyle \text{sq mean implies three proportional}$\par
\medskip
\noindent
\begin{tabular}{@{} >{\raggedright\arraybackslash}p{0.06\textwidth} >{\raggedright\arraybackslash}p{0.47\textwidth} >{\raggedleft\arraybackslash}p{0.41\textwidth} @{}}
\textbf{\#} & \textbf{Proof} & \textbf{Mathematics} \\
\midrule
\textbf{1.} & We prove that proposition 18: rectangle on extremes equal to square on mean implies three lines proportional for Euclid Book VI the Euclidean plane. & \\[0.45em]
\textbf{2.} & Let a = first. & \\[0.45em]
\textbf{3.} & Let b = second. & \\[0.45em]
\textbf{4.} & Let c = third. & \\[0.45em]
\textbf{5.} & From step 2, step 3, and step 4, we can deduce that ratio(a, b) equals ratio(b, c). & $\displaystyle ratio(a, b) = ratio(b, c)$ \\[0.45em]
\textbf{6.} & We invoke the derived fact governing Euclid Book VI the Euclidean plane: proposition 17: three proportional lines give rectangle on extremes equal to square on mean, for Euclid Book VI on the Euclidean plane. & \\[0.45em]
\textbf{7.} & From step 6, this implies sq mean implies three proportional. Hence proven. & $\displaystyle \text{sq mean implies three proportional}$ \\[0.45em]
\bottomrule
\end{tabular}
\label{thm:Geometry-Euclid Book VI-proposition_18_rectangle_on_extremes_equal_to_square_on_mean_implies_three_lines_proportional}
\par\medskip\hrule
\end{document}
